% Lösungen Einheit 3 von 5 – Kumulierte Wahrscheinlichkeiten (zusammenbau v0.1)
\einheitenkopf{Einheit 3 von 5 · Kumulierte Wahrscheinlichkeiten}

\erg{15}{a) $P(X \le 4)$ – weniger als fünf heißt höchstens vier \quad b) Grenze 6: $P(X \le 6)$ \quad c) $P(X \le 7) \approx 0{,}3061$ (X: Anzahl, n = 25) \quad d) $1 - P(X \le 14) \approx 0{,}3552$ (mindestens 15; X: Anzahl) \quad e) $P(A) = P(X \le 6) \approx 0{,}3643$; $P(B) = P(X \le 11) - P(X \le 8) \approx 0{,}2674$ \quad f) $P(A) = P(X = 3) \approx 0{,}1875$; 10\,\% von 60 sind 6: $P(B) = 1 - P(X \le 5) \approx 0{,}2420$ \quad g) $60 - H \ge 2H$, also $H \le 20$: $P(H \le 20) \approx 0{,}4516$ \quad h) $1 - (P(X \le 198) - P(X \le 161)) \approx 0{,}1414$ (Erwartungswert 180, mehr als 18 Abweichung: X < 162 oder X > 198) \quad i) Wahrscheinlichkeit, dass mindestens zwei der kontrollierten Räder kein funktionierendes Licht haben (eins minus: keines oder genau eines). \quad j) Extrapreis: vorn genau vier Kreise, $(6 \text{ über } 4) \cdot 0{,}4^4 \cdot 0{,}6^2$; dann braucht der Gewinn mindestens einen Kreis unter den letzten zwei, $1 - 0{,}6^2$. $P \approx 0{,}0885$}
\erg{16}{a) 52\,\% von 400 sind 208, also höchstens 150 Zahl in den restlichen 300 Würfen: $P(Y \le 150) \approx 0{,}5230$ (Y: Zahl im Rest)}
\erg{17}{a) $(1 - P(A \le 5)) \cdot (1 - P(B \le 5)) \approx 0{,}2387$ (A, B: Treffer von Anna und Ben)}
\erg{18}{a) Fehler: Paul liest „mehr als 7“ als „mindestens 7“. Das Gegenereignis ist „höchstens 7“. Richtig: $1 - P(X \le 7) \approx 0{,}3990$ \quad b) Das Gegenteil von „mindestens k“ sind alle Anzahlen darunter, also bis k − 1; die Anzahl k selbst gehört zu „mindestens k“ und darf nicht abgezogen werden. \quad c) Die Zimmer reichen bei mindestens 5 Stornos: $1 - P(X \le 4) \approx 0{,}8757$ (X: Anzahl der Stornos) \quad d) $\Sigma_{k=0}^{3}\, (25 \text{ über } k) \cdot 0{,}15^k \cdot 0{,}85^{25-k} \approx 0{,}4711$}
